% TOP_PROBLEM: {"id": 30006453, "title": "Maximum Normalized Mahler Measure of Littlewood Polynomials", "category_id": 1, "category": {"id": 1, "name": "number_theory", "display_name": "Number Theory", "description": "Properties of integers, prime numbers, Diophantine equations.", "slug": "number-theory", "order_index": 1, "created_at": "2026-07-31T15:26:25.670Z"}, "status": "partially_solved"}
% FIRST_POSTED: 2026-10-06
% ATTEMPT_STATUS: unresolved
% OPENAI_CLAIM: {"schema_version":1,"source_repo":"https://github.com/openai/math","source_commit":"adc7f1241b42e322a6451854ab7e4b4c146bf78a","release_date":"2026-10-06","match":"partial","resolution":"partial","independently_reviewed":false,"second_pass":"n/a","families":[{"family":"076","title":"Real ultraflat Littlewood polynomials and unbounded binary merit factors","subject":"Real and complex analysis","manuscripts":[{"title":"Asymptotically minimal maxima of real Littlewood polynomials","pdf_path":"preprints/Asymptotically-minimal-maxima-of-real-Littlewood-polynomials-September-23-2026/paper.pdf","readme_path":"preprints/Asymptotically-minimal-maxima-of-real-Littlewood-polynomials-September-23-2026/README.md","date":"2026-09-23","theorem_ref":"Theorem/corollary (label thm:main), preprints/Asymptotically-minimal-maxima-of-real-Littlewood-polynomials-September-23-2026/build/introduction.tex"},{"title":"Nearly minimal maxima and positive minima of Littlewood polynomials","pdf_path":"preprints/Nearly-minimal-maxima-and-positive-minima-of-Littlewood-polynomials-October-5-2026/littlewood-lower-envelope.pdf","readme_path":"preprints/Nearly-minimal-maxima-and-positive-minima-of-Littlewood-polynomials-October-5-2026/README.md","date":"2026-10-05","theorem_ref":"Theorem/corollary (label thm:main), preprints/Nearly-minimal-maxima-and-positive-minima-of-Littlewood-polynomials-October-5-2026/build/sections/introduction.tex"},{"title":"Ultraflat real Littlewood polynomials","pdf_path":"preprints/Ultraflat-real-Littlewood-polynomials-October-5-2026/ultraflat-real-littlewood-polynomials.pdf","readme_path":"preprints/Ultraflat-real-Littlewood-polynomials-October-5-2026/README.md","date":"2026-10-05","theorem_ref":"Theorem/corollary (label thm:main), preprints/Ultraflat-real-Littlewood-polynomials-October-5-2026/build/sections/introduction.tex"}],"lean":{"doc_path":"lean/docs/076.md","comparators":["lean/ComparatorChallenges/AsymptoticallyMinimalLittlewood.json","lean/ComparatorChallenges/AsymptoticallyMinimalLittlewood.lean","lean/ComparatorChallenges/LittlewoodFiniteFlatness.json","lean/ComparatorChallenges/LittlewoodFiniteFlatness.lean"],"declarations":["OAI.AsymptoticallyMinimalLittlewood.main"],"covers_main_theorem":false},"reasoning_trace":null}]}
\documentclass[11pt]{article}
\usepackage[margin=25mm]{geometry}
\usepackage{fontspec}
\setmainfont{Latin Modern Roman}
\usepackage{amsmath,amssymb,mathtools}
\usepackage{unicode-math}
\setmathfont{Latin Modern Math}
\usepackage{enumitem}
\usepackage{xurl,hyperref,bookmark}
\hypersetup{colorlinks=true,urlcolor=blue}
\setcounter{secnumdepth}{0}
\setlength{\parindent}{0pt}
\setlength{\parskip}{5pt}
\setlength{\emergencystretch}{3em}
\begin{document}
\section*{Maximum Normalized Mahler Measure of Littlewood Polynomials}

\subsection{OpenAI's claim}
\textit{Nearly optimal Littlewood maxima}: The optimal supremum norm of a real sign polynomial of length $N$ is claimed to approach $\sqrt N$ relatively, with $N$ running through every integer; the resulting binary sequences have merit factors tending to infinity.

\textit{Littlewood polynomials with a lower margin}: The paper claims sign polynomials at every sufficiently large length with maximum at most $(1+\eta)\sqrt N$ and minimum at least $\sqrt N/16$, for any fixed positive $\eta$; both inequalities concern the same polynomial.

\textit{Uniformly flat real sign polynomials}: OpenAI claims that real sign polynomials can have normalized modulus uniformly as close to one as desired once the integer length is large enough; choices at different lengths are independent of one another.

\subsection{Outline of the argument}
\textit{Nearly optimal Littlewood maxima}: The construction starts with bounded real coefficients obtained by sampling quadratic-phase waves on an auxiliary torus. Interval packing separates the relevant frequency ranges and keeps almost all of the mean-square mass. A discrepancy rounding step changes the coefficients to signs while keeping the circle supremum error small compared with $\sqrt N$; Parseval’s identity then supplies the lower comparison.

\textit{Littlewood polynomials with a lower margin}: The relaxed-coefficient construction for small maxima is augmented on arcs where its leading amplitude is too small. Carefully matched quadratic phases create a positive lower margin without introducing excessive Fourier coefficients outside the permitted degree range. A rounding estimate expressed in terms of the remaining coefficient defect preserves the margin while turning the coefficients into signs.

\textit{Uniformly flat real sign polynomials}: A continuous circle function is built with nearly constant modulus, real coefficients of nearly sign size in the retained frequency range, and a small Fourier tail. Auxiliary-torus waves and phase corrections coordinate the upper and lower modulus bounds. Parseval controls the aggregate coefficient defect, allowing discrepancy rounding to preserve uniform flatness.

\subsection{Scope relative to this problem}
\textit{Nearly optimal Littlewood maxima}: The constructed maxima approach sqrt(N) and the normalized moduli converge in every finite positive L\textasciicircum{}p, but no uniform lower bound or logarithmic uniform integrability is supplied. Mahler measure is a logarithmic mean, so its extremal normalized limit is not established by this upper-bound theorem alone.

\textit{Littlewood polynomials with a lower margin}: The target limsup is the largest real-sign Mahler measure divided by sqrt(N). The polynomial has upper modulus (1+o(1))sqrt(N), L2 norm exactly sqrt(N), and a uniform positive lower multiple of sqrt(N). The upper bound and L2 identity force modulus/sqrt(N) to approach one in measure; bounded logarithms on the fixed-positive-lower-bound interval then force the mean log modulus to zero. Jensen gives the universal upper bound one, so the limsup equals one.

\textit{Uniformly flat real sign polynomials}: The target limsup is the largest real-sign Mahler measure divided by sqrt(N). The polynomial has upper modulus (1+o(1))sqrt(N), L2 norm exactly sqrt(N), and a uniform positive lower multiple of sqrt(N). The upper bound and L2 identity force modulus/sqrt(N) to approach one in measure; bounded logarithms on the fixed-positive-lower-bound interval then force the mean log modulus to zero. Jensen gives the universal upper bound one, so the limsup equals one.

This page records partial support. Complete resolution of the catalogue question is not established by this record.

\subsection{Formal verification}
Family 076 describes Lean coverage for near-minimal upper maxima and convergence of finite-exponent modulus means; the later uniform lower-envelope conclusions are separate.

\textit{Nearly optimal Littlewood maxima}: The manuscript links \texttt{AsymptoticallyMinimalLittlewood.lean} (declaration \texttt{OAI.AsymptoticallyMinimalLittlewood.main}), \texttt{LittlewoodFiniteFlatness.lean}.

The documented selection includes this manuscript’s main assertion; that assertion retains the restrictions stated above.

\textit{Littlewood polynomials with a lower margin}: No manuscript-specific Comparator is listed in the pinned inventory.

The documented selection does not cover the entire manuscript assertion relevant here.

\textit{Uniformly flat real sign polynomials}: No manuscript-specific Comparator is listed in the pinned inventory.

The documented selection does not cover the entire manuscript assertion relevant here.

This site has not run the Lean development. These links describe OpenAI’s selected formal statements and should be read with their stated scope.

\subsection{Status caveat}
OpenAI reports that an internal model produced the manuscripts. Their arguments have not been independently verified as proofs by this site. Any portion outside the documented formal scope may contain mistakes, and the catalogue label reports the announced mathematical claim.

\subsection{Sources}
\begin{itemize}
\item OpenAI, \emph{Asymptotically minimal maxima of real Littlewood polynomials}, 2026-09-23; citation key \texttt{OAI:Asymptotically-minimal-maxima-of-real-Littlewood-polynomials-September-23-2026}. \href{https://github.com/openai/math/blob/adc7f1241b42e322a6451854ab7e4b4c146bf78a/preprints/Asymptotically-minimal-maxima-of-real-Littlewood-polynomials-September-23-2026/paper.pdf}{PDF (GitHub)}; \href{https://github.com/openai/math/raw/adc7f1241b42e322a6451854ab7e4b4c146bf78a/preprints/Asymptotically-minimal-maxima-of-real-Littlewood-polynomials-September-23-2026/paper.pdf}{PDF (direct)}; \href{https://github.com/openai/math/blob/adc7f1241b42e322a6451854ab7e4b4c146bf78a/preprints/Asymptotically-minimal-maxima-of-real-Littlewood-polynomials-September-23-2026/README.md}{Manuscript page and citation}; \href{https://github.com/openai/math/blob/main/preprints/Asymptotically-minimal-maxima-of-real-Littlewood-polynomials-September-23-2026/paper.pdf}{Latest version}.
\item OpenAI, \emph{Nearly minimal maxima and positive minima of Littlewood polynomials}, 2026-10-05; citation key \texttt{OAI:Nearly-minimal-maxima-and-positive-minima-of-Littlewood-polynomials-October-5-2026}. \href{https://github.com/openai/math/blob/adc7f1241b42e322a6451854ab7e4b4c146bf78a/preprints/Nearly-minimal-maxima-and-positive-minima-of-Littlewood-polynomials-October-5-2026/littlewood-lower-envelope.pdf}{PDF (GitHub)}; \href{https://github.com/openai/math/raw/adc7f1241b42e322a6451854ab7e4b4c146bf78a/preprints/Nearly-minimal-maxima-and-positive-minima-of-Littlewood-polynomials-October-5-2026/littlewood-lower-envelope.pdf}{PDF (direct)}; \href{https://github.com/openai/math/blob/adc7f1241b42e322a6451854ab7e4b4c146bf78a/preprints/Nearly-minimal-maxima-and-positive-minima-of-Littlewood-polynomials-October-5-2026/README.md}{Manuscript page and citation}; \href{https://github.com/openai/math/blob/main/preprints/Nearly-minimal-maxima-and-positive-minima-of-Littlewood-polynomials-October-5-2026/littlewood-lower-envelope.pdf}{Latest version}.
\item OpenAI, \emph{Ultraflat real Littlewood polynomials}, 2026-10-05; citation key \texttt{OAI:Ultraflat-real-Littlewood-polynomials-October-5-2026}. \href{https://github.com/openai/math/blob/adc7f1241b42e322a6451854ab7e4b4c146bf78a/preprints/Ultraflat-real-Littlewood-polynomials-October-5-2026/ultraflat-real-littlewood-polynomials.pdf}{PDF (GitHub)}; \href{https://github.com/openai/math/raw/adc7f1241b42e322a6451854ab7e4b4c146bf78a/preprints/Ultraflat-real-Littlewood-polynomials-October-5-2026/ultraflat-real-littlewood-polynomials.pdf}{PDF (direct)}; \href{https://github.com/openai/math/blob/adc7f1241b42e322a6451854ab7e4b4c146bf78a/preprints/Ultraflat-real-Littlewood-polynomials-October-5-2026/README.md}{Manuscript page and citation}; \href{https://github.com/openai/math/blob/main/preprints/Ultraflat-real-Littlewood-polynomials-October-5-2026/ultraflat-real-littlewood-polynomials.pdf}{Latest version}.
\item Family 076: \href{https://github.com/openai/math/blob/adc7f1241b42e322a6451854ab7e4b4c146bf78a/lean/docs/076.md}{Lean scope}.
\item \href{https://github.com/openai/math/blob/adc7f1241b42e322a6451854ab7e4b4c146bf78a/lean/ComparatorChallenges/AsymptoticallyMinimalLittlewood.json}{Comparator AsymptoticallyMinimalLittlewood.json}.
\item \href{https://github.com/openai/math/blob/adc7f1241b42e322a6451854ab7e4b4c146bf78a/lean/ComparatorChallenges/AsymptoticallyMinimalLittlewood.lean}{Comparator AsymptoticallyMinimalLittlewood.lean}.
\item \href{https://github.com/openai/math/blob/adc7f1241b42e322a6451854ab7e4b4c146bf78a/lean/ComparatorChallenges/LittlewoodFiniteFlatness.json}{Comparator LittlewoodFiniteFlatness.json}.
\item \href{https://github.com/openai/math/blob/adc7f1241b42e322a6451854ab7e4b4c146bf78a/lean/ComparatorChallenges/LittlewoodFiniteFlatness.lean}{Comparator LittlewoodFiniteFlatness.lean}.
\item \href{https://github.com/openai/math/blob/adc7f1241b42e322a6451854ab7e4b4c146bf78a/CONTENTS.md}{OpenAI catalogue entry}.
\end{itemize}
\end{document}
